\subsection{\texorpdfstring{Relations between \(\mathcal{L}(\mathbf{E};\mathbf{F})\) and \(\mathcal{L}(\hat{\mathbf{E}};\mathbf{F})\)}{Relations between \textbackslash mathcal\{L\}(\textbackslash mathbf\{E\};\textbackslash mathbf\{F\}) and \textbackslash mathcal\{L\}(\textbackslash hat\{\textbackslash mathbf\{E\}\};\textbackslash mathbf\{F\})}}

Let E and F the two Hausdorff locally convex spaces, and suppose F is complete; let \(\hat{E}\) be the completion of E. Since every continuous linear mapping u from E into F extends uniquely to a continuous linear mapping \(\overline{u}\) from \(\hat{E}\) into F, we can identify the spaces \(\mathcal{L}(E;F)\) and \(\mathcal{L}(\hat{E};F)\) by the mapping \(u\mapsto\overline{u}\) . In addition, let S be a family of bounded subsets of E; the S-topology on \(\mathcal{L}(E;F)\) coincides with the S-topology on \(\mathcal{L}(\hat{E};F)\) and also with the \(\hat{S}\) -topology, where \(\hat{S}\) denotes the family of closures in \(\hat{E}\) of sets of S.

For example, if E is normed, the topology of bounded convergence on \(\mathcal{L}(\mathrm{E};\mathrm{F})\) is identical with the topology of bounded convergence on \(\mathcal{L}(\hat{\mathrm{E}};\mathrm{F})\) : for, every bounded subset of \(\hat{E}\) is contained in the closure of a bounded subset of E. Since the unit ball of \(\hat{E}\) is the closure of the unit ball of E, it follows from formula (3) (III, p.~14) that if F is a Banach space, the map \(u\mapsto\overline{u}\) is an isometry from \(\mathcal{L}(\mathrm{E};\mathrm{F})\) onto \(\mathcal{L}(\hat{\mathrm{E}};\mathrm{F})\) .

We observe that if E is not a normed space, then there may exist bounded subsets of \(\hat{E}\) which are not contained in the closure of any bounded subset of E (for example, if E is the weak dual of an infinite dimensional Banach space); however, this is so if E is metrizable and satisfies the first axiom of countability (III, p.~39, exerc. 16).

